
\chapterwithtoc{Conclusion générale}
\label{chap:conclusion}

Cet essai partait d'un constat : le goulot d'étranglement d'un centre
d'opérations de sécurité n'est pas la détection, mais la qualification de ce qui
est détecté. Il a conçu et mesuré un écosystème ancrant un modèle de langage sur
MITRE ATT\&CK par le \emph{Model Context Protocol}.

\textbf{Les modèles seuls ne suffisent pas, et l'écart se creuse côté
industriel.} Sur 42 échantillons produits en laboratoire, hors contamination, dix
modèles courants ont été évalués en 1\,260 appels. Aucun ne dépasse 0,603 de F1
technique en base stricte et les cinq premiers se tiennent en moins de 0,04 : il
n'y a pas de vainqueur net. La ventilation par matrice est plus instructive que
le classement --- le meilleur passe de 0,679 en Enterprise à 0,452 en ICS. Les
vingt-trois identifiants fabriqués du panel sont tous des tactiques, et tous sur
des journaux industriels : l'angle mort signalé par la revue est mesuré, et il
est d'abord un défaut de repérage entre univers.

\textbf{L'ancrage corrige ce que le modèle ne peut pas corriger seul.} À corpus
et conditions identiques, le modèle retenu progresse de 0,563 à 0,706 en
technique et de 0,587 à 0,802 en tactique. Le gain n'est pas uniforme : la
matrice industrielle double presque, de 0,262 à 0,524, ce qui écarte
l'hypothèse d'un effet de consigne. Surtout, le taux d'identifiants obsolètes ne
diminue pas, il disparaît --- une classe entière d'erreurs supprimée, qu'aucun
raisonnement ne rattrape puisque rien dans le journal ne signale que le
référentiel a bougé depuis l'entraînement.

\textbf{Le mérite revient à la conception du serveur}, non au simple accès à une
source : une version nommée et jointe à chaque réponse, la résolution des
révocations, une granularité d'outils bornée, une trace d'usage vérifiable. Le
risque, lui, est déplacé et non supprimé : le serveur devient un canal
d'instruction autant qu'une source de faits, et l'audit de cette surface reste à
mener.

Trois limites bornent ces résultats : un seul modèle évalué avec ancrage, un
corpus de 42 échantillons, et un effort de raisonnement laissé aux valeurs par
défaut faute d'échelle commune entre éditeurs. Reste l'enseignement principal :
les modèles les plus capables sont les plus exposés, parce que l'attaque exploite
l'aptitude même qui fait leur valeur. La réponse n'est pas un modèle plus grand,
mais un cadre qui ancre, borne et journalise.
